% Drawing 156-7/04-2 -- folder 156-7, batch 04, archive page 69 (PDF page 70).
% Transcription: à gauche, sous la formule, trois petits croquis : devant « b) », un segment marqué $X$, son extrémité gauche un gros point que désigne une petite flèche partie de $Y$ ; devant « c) », deux segments séparés marqués $X$ et $Y$, puis « ou », et au-dessous deux segments $X$, $Y$ bout à bout, un gros point à leur extrémité commune.
% Kind: segments and points on a line (cases b and c of compatibility). Method 2 (schematic redraw from the facsimile). Status: ok.
% Notes: Three small sketches at the lower left of the page, illustrating the cases (b) and (c) of the text. After « b) »: a segment labelled X, a heavy point at its left end, to which a short arrow points down from the letter Y; the right end of the segment is thickened. After « c) »: two separate segments labelled X and Y (X with a heavy point at its left end and a tick at its right end, Y with a tick at each end), then « ou » to the right, and below them two segments X, Y end to end, a tick at each outer end and a heavy point at the common end. Labels as on the page. Nothing struck. Not encoded as data (segment-config): the three rows are separate configurations with a text label each; redrawn. Redrawn 2026-11-10 (issue 128): to check again.
% Coordinates: 1 unit = 1 mm = 1/100 of the crop width (bbox in page fractions: [0.11, 0.76, 0.37, 0.84]).
% Redrawn by Claude (claude-fable-5-2), issue #128, 2026-11-10; not checked by a human; not compiled (no TeX in the session).
% Rights: reproduces part of the work -- behind RENDER_TRANSCRIPTS, never published by itself.
\documentclass[tikz,border=3pt]{standalone}
\usepackage{amsmath,amssymb}
\usetikzlibrary{patterns,arrows.meta}
\tikzset{pencil/.style={black!45}, struck/.style={black!35}, mpt/.style={circle,draw,inner sep=1.3pt}, pt/.style={circle,fill,inner sep=0.9pt}, opt/.style={circle,draw,fill=white,inner sep=1.1pt}}
\begin{document}
\begin{tikzpicture}[x=1mm,y=1mm,line cap=round,line join=round,font=\small]
\node[inner sep=1pt] at (10.15,-18.62) {\small b)};
\draw[thick] (25.64,-18.02) -- (46.54,-16.75);
\node[pt, inner sep=1.4pt] at (25.64,-18.02) {};
\draw[very thick] (45.27,-16.84) -- (46.54,-16.75);
\node[inner sep=1pt] at (24.12,-9.73) {\small $Y$};
\draw[-{Stealth[length=1.2mm]}] (25.47,-12.27) -- (25.64,-16.08);
\node[inner sep=1pt] at (36.38,-12.27) {\small $X$};
\node[inner sep=1pt] at (10.15,-29.62) {\small c)};
\draw[thick] (22.85,-29.19) -- (39.77,-28.35);
\node[pt, inner sep=1.2pt] at (22.85,-29.19) {};
\draw[thick] (39.77,-27.33) -- (39.77,-29.36);
\node[inner sep=1pt] at (31.31,-24.96) {\small $X$};
\draw[thick] (52.46,-27.75) -- (68.96,-27.42);
\draw[thick] (52.46,-26.48) -- (52.46,-29.02);
\draw[thick] (68.96,-26.15) -- (68.96,-28.68);
\node[inner sep=1pt] at (63.46,-23.69) {\small $Y$};
\node[inner sep=1pt] at (92,-27.5) {\small ou};
\draw[thick] (22.42,-37.91) -- (39.77,-37.65) -- (57.54,-37.4);
\draw[thick] (22.42,-36.72) -- (22.42,-39.09);
\node[pt, inner sep=1.5pt] at (39.77,-37.65) {};
\draw[thick] (57.54,-36.22) -- (57.54,-38.58);
\node[inner sep=1pt] at (30.88,-34.27) {\small $X$};
\node[inner sep=1pt] at (49.5,-33.85) {\small $Y$};
\end{tikzpicture}
\end{document}
